\section{RPC and streaming}\label{sec:session_rpc}

The session layer builds remote procedure calls on top of ordinary pub-sub
messages rather than introducing a separate service-ID abstraction. An RPC
request is a normal reliable publication whose recipients may emit one or more
unicast \emph{responses} back to the publisher. The number of responses per
request is unrestricted; streaming responses are therefore a zero-cost
extension of single-response RPC.

\subsection{Breadcrumbs}

When a subscriber receives a message, the transport supplies the unicast
metadata needed to reach the originating node. The subscriber captures this
information in a \emph{breadcrumb} consisting of:

\begin{itemize}
    \item \textit{remote\_id}: the network identifier of the original
    publisher;
    \item \textit{topic\_hash}: the hash of the topic the original message
    was delivered on;
    \item \textit{message\_tag}: the tag of the original message;
    \item \textit{seqno}: a 48-bit counter, initialized to zero, that is
    incremented by the subscriber once for each response it sends in reply
    to this message;
    \item \textit{unicast\_ctx}: the transport's unicast metadata for the
    original publisher.
\end{itemize}

A breadcrumb is sufficient to emit any number of responses to the original
publisher at a later time, without keeping the original message alive.

\subsection{Responses}

A response is carried by a \texttt{rsp\_be} or \texttt{rsp\_rel} transfer
(section~\ref{sec:session_header_rsp}). The response header carries:

\begin{itemize}
    \item \textit{topic\_hash} and \textit{message\_tag}, copied from the
    breadcrumb, which together identify the request;
    \item \textit{seqno}, the 48-bit response sequence number;
    \item \textit{response\_tag}, an arbitrary byte chosen by the responder
    for acknowledgement correlation (reliable responses only).
\end{itemize}

Streaming is achieved simply by sending multiple responses to the same
request, each with an incremented \textit{seqno}. The publisher observes a
monotonic stream of responses per remote responder.

\subsection{Reliable responses}

A responder that wishes to detect whether the original publisher is still
listening shall send its responses reliably (\texttt{rsp\_rel}). The
publisher acknowledges a reliable response with either a \texttt{rsp\_ack}
or a \texttt{rsp\_nack}:

\begin{description}
    \item[\texttt{rsp\_ack}] The publisher's application is still accepting
    responses for this request. The responder may continue to stream.

    \item[\texttt{rsp\_nack}] The publisher's application has closed the
    request and will discard any further responses. The responder should
    cease streaming promptly.
\end{description}

\subsection{Client-side deduplication and ordering}

The publisher tracks, for each remote responder, a small bitmap of recently
seen \textit{seqno} values and rejects duplicates within the bitmap window.
Responses whose \textit{seqno} falls outside the window because they are too
old are rejected.

Ordering is not reconstructed by default: a response with a larger
\textit{seqno} may be delivered before one with a smaller \textit{seqno}
on the same stream. An application that requires strict ordering within a
stream may implement a reorder buffer keyed on \textit{seqno}; pattern-based
subscribers may opt into such a buffer via the subscriber's ordering
configuration.

\subsection{Liveness monitoring}

The publisher may attach a \emph{liveness timeout} to a pending request.
If no response has been received from \emph{any} remote for longer than the
timeout, the publisher's observable state is moved into a \textsc{liveness}
warning that alerts the application and is automatically cleared when the
next response arrives. Liveness state is independent of the response queue:
a pending response that has not yet been consumed does not inhibit the
liveness warning, and clearing the warning does not drop the pending
response.

\subsection{Zombie requests}

After the publisher's application has closed a request, the session layer
may retain a small amount of state for the request briefly in order to
correctly handle duplicates of reliable responses that arrive after the
application has ceased to listen. A duplicate that arrives in this zombie
window is acknowledged with the same positive or negative acknowledgement
as the first time, without involving the application. The zombie window
expires with the session lifetime (section~\ref{sec:session_parameters}).

\begin{remark}
    The zombie mechanism exists purely to absorb benign duplicates. It is
    not a mechanism for deferring application cleanup or for re-dispatching
    responses to a replacement handler.
\end{remark}

\section{Pattern subscriptions}\label{sec:session_patterns}

Pattern subscriptions are indistinguishable from verbatim subscriptions on
the wire. The distinction exists only in how the local node decides which
topics to join.

\subsection{Discovery}

A node with a local pattern subscription discovers matching topics in two
complementary ways:

\begin{description}
    \item[Passive discovery] A received gossip carries the topic name (see
    section~\ref{sec:session_header_gossip}). If the received name matches a
    local pattern, the node creates an \emph{implicit topic} for it and
    attaches a subscriber. Passive discovery is automatic and requires no
    additional traffic.

    \item[Active discovery] A node that has just installed a new pattern
    subscription may broadcast a scout message
    (section~\ref{sec:session_scout}) with the pattern as the payload. Other
    nodes respond with gossip messages for any of their topics that match,
    which drives passive discovery.
\end{description}

\subsection{Substitutions}

When a pattern matches a name, the pattern's wildcards $(\texttt{*},\ \texttt{>})$
are paired with the segments they matched. The result is an ordered list of
\emph{substitutions}, each recording the matched substring together with its
ordinal position in the pattern.

\begin{remark}
    For example, the pattern \texttt{ins/*/data/>} matching the name
    \texttt{ins/0/data/foo/456} yields three substitutions:
    \verb|"0"|, \verb|"foo"|, and \verb|"456"|, at ordinals $0$, $1$, and $1$
    respectively (the second and third substitutions share an ordinal
    because they are both consumed by the trailing \texttt{>} token).
\end{remark}

The substitutions are computed once, when the pattern first matches the
topic, and are made available to the receiving application alongside each
arriving message. No string re-matching is performed per message.

\subsection{Implicit topic lifecycle}

An implicit topic behaves like any other topic for arbitration and message
delivery, with two differences:

\begin{itemize}
    \item It does not emit gossip of its own. Its CRDT state is updated by
    incoming gossip, by inline gossip from incoming application messages, or
    by scout-triggered passive discovery.
    \item It is automatically retired after the retention window elapses
    without any activity (no messages, no gossips).
\end{itemize}

The retention window should be much longer than the worst-case periodic
gossip interval so that a silent but alive topic is not prematurely
retired. A recommended default is ten minutes; see
section~\ref{sec:session_parameters}.

\subsection{Publisher invariant}

A publisher's topic name shall not contain any pattern tokens. A resolved
publisher name that contains \texttt{*} or \texttt{>} is invalid and shall
be rejected by the session layer at publisher creation time. This invariant
ensures that every topic on the wire has a deterministic, unambiguous name.
